\documentclass{article}

\usepackage{fancyhdr}
\usepackage{extramarks}
\usepackage{amsmath}
\usepackage{amsthm}
\usepackage{amsfonts}
\usepackage{tikz}
\usepackage[plain]{algorithm}
\usepackage{algpseudocode}
\usepackage{fullpage, url, hyperref}
\newcommand{\dotcup}{\ensuremath{\mathaccent\cdot\cup}}

\theoremstyle{plain}
\newtheorem{theorem}{Theorem}
\newtheorem{lemma}{Lemma}
\newtheorem{definition}{Definition}

\newcommand{\I}{\mathcal{I}}
\newcommand{\F}{\mathbb{F}}
\newcommand{\E}{\mathbb{E}}
\newcommand{\R}{\mathbb{R}}
\newcommand{\tr}{\mathrm{tr}}
\renewcommand{\P}{\mathbb{P}}
\usetikzlibrary{automata,positioning}

%
% Basic Document Settings
%

\topmargin=-0.45in
\evensidemargin=-.50in
\oddsidemargin=-.5in
\textwidth=7in
\textheight=9.5in
\headsep=0.25in

\linespread{1.1}

\pagestyle{fancy}
\lhead{\hmwkAuthorName}
\chead{\hmwkClass\ (\hmwkClassInstructor\ \hmwkClassTime): \hmwkTitle}
\rhead{\firstxmark}
\lfoot{\lastxmark}
\cfoot{\thepage}

\newcommand{\w}[2]{w_{#1}^{(#2)}}
\newcommand{\loss}[2]{\vec{\ell}_{#2}^{(#1)}}
\newcommand{\losst}[1]{\vec{\ell}^{(#1)}}
  \newcommand{\pt}[1]{\vec{p}^{(#1)}}
\renewcommand\headrulewidth{0.4pt}
\renewcommand\footrulewidth{0.4pt}

\setlength\parindent{0pt}


% 
% Create Problem Sections
%


\newcommand{\enterProblemHeader}[1]{
    \nobreak\extramarks{}{}\nobreak{}
    \nobreak\extramarks{}\nobreak{}
}

\newcommand{\exitProblemHeader}[1]{
    \nobreak\extramarks{}\nobreak{}
    \stepcounter{#1}
    \nobreak\extramarks{}\nobreak{}
  }

  
\newcommand{\enterExerciseHeader}[1]{
    \nobreak\extramarks{}{}\nobreak{}
    \nobreak\extramarks{}\nobreak{}
}

\newcommand{\exitExerciseHeader}[1]{
    \nobreak\extramarks{}\nobreak{}
    \stepcounter{#1}
    \nobreak\extramarks{}\nobreak{}
}

\setcounter{secnumdepth}{0}
\newcounter{partCounter}
\newcounter{homeworkProblemCounter}
\setcounter{homeworkProblemCounter}{1}

\newcounter{homeworkExerciseCounter}
\setcounter{homeworkExerciseCounter}{1}
\nobreak\extramarks{}\nobreak{}
\nobreak\extramarks{}\nobreak{}

%
% Homework Problem Environment
%
% This environment takes an optional argument. When given, it will adjust the
% problem counter. This is useful for when the problems given for your
% assignment aren't sequential. See the last 3 problems of this template for an
% example.
%
\newenvironment{homeworkProblem}[1][-1]{
%    \ifnum#1>0
   %     \setcounter{homeworkProblemCounter}{#1}
   % \fi
    \section{Problem \arabic{homeworkProblemCounter}}
    \setcounter{partCounter}{1}
    \enterProblemHeader{homeworkProblemCounter}
}{
    \exitProblemHeader{homeworkProblemCounter}
}

  \newenvironment{homeworkExercise}[1][-1]{
    \ifnum#1>0
        \setcounter{homeworkExerciseCounter}{#1}
    \fi
    \section{Exercise \arabic{homeworkExerciseCounter}}
    \setcounter{partCounter}{1}
    \enterExerciseHeader{homeworkExerciseCounter}
}{
    \exitExerciseHeader{homeworkExerciseCounter}
}

%
% Homework Details
%   - Title
%   - Due date
%   - Class
%   - Section/Time
%   - Instructor
%   - Author
%

\newcommand{\hmwkTitle}{Homework\ \#2}
\newcommand{\hmwkAssignedDate}{October 24}
\newcommand{\hmwkDueDate}{Novemer 5, midnight}
\newcommand{\hmwkClass}{Foundations of Fairness in Machine Learning, Fall 2018}
\newcommand{\hmwkClassTime}{}%Section A}
\newcommand{\hmwkClassInstructor}{Professor Jamie Morgenstern}
\newcommand{\hmwkAuthorName}{}%\textbf{Josh Davis} \and \textbf{Davis Josh}}

%
% Title Page
%

\title{    \vspace{2in}
    \textmd{\textbf{\hmwkClass:\ \hmwkTitle}}\\
    \normalsize\vspace{0.1in}\small{Due\ on\ \hmwkDueDate\ at 3:10pm}\\
    \vspace{0.1in}\large{\textit{\hmwkClassInstructor\ \hmwkClassTime}}
    \vspace{3in}
}

\author{\hmwkAuthorName}
\date{}

\renewcommand{\part}[1]{\textbf{\large Part \Alph{partCounter}}\stepcounter{partCounter}\\}

%
% Various Helper Commands
%

% Useful for algorithms
\newcommand{\alg}[1]{\textsc{\bfseries \footnotesize #1}}

% For derivatives
\newcommand{\deriv}[1]{\frac{\mathrm{d}}{\mathrm{d}x} (#1)}

% For partial derivatives
\newcommand{\pderiv}[2]{\frac{\partial}{\partial #1} (#2)}

% Integral dx
\newcommand{\dx}{\mathrm{d}x}

% Alias for the Solution section header
\newcommand{\solution}{\textbf{\large Solution}}

% Probability commands: Expectation, Variance, Covariance, Bias
\newcommand{\Var}{\mathrm{Var}}
\newcommand{\Cov}{\mathrm{Cov}}
\newcommand{\Bias}{\mathrm{Bias}}
\usepackage{enumitem}
\newcommand\litem[1]{\item{\bfseries#1.\space}}

\newcommand{\Z}{\mathbb{Z}}

\newcommand{\poly}{\textrm{poly}}
\begin{document}

% \maketitle

\pagebreak

{\bf Homework Out: \hmwkAssignedDate}\\
{\bf Due Date:  \hmwkDueDate}\\
{\bf Reminder: }

To submit your homework, please go to
\url{https://classroom.github.com/a/7rslWmMd}, accept the assignment,
and submit your LaTex, PDF, and any code you use for the
assignment. Please name your files ``hw2-USERNAME-writeup.{tex,pdf}''
``hw2-USERNAME-code.{appropriate file type}''. You will need a github
account, and to add, commit, and push your homework.

Please cite all sources you use, and people you work with. The
expectation is that you try and solve these problems yourself, rather
than looking online explicitly for answers. Submissions due at 23:00
of the due date.

{\bf You may use $O$-notation unless explicitly noted somewhere in the homework.}


\newcommand{\X}{\mathcal{X}}
\newcommand{\Y}{\mathcal{Y}}
\newcommand{\D}{\mathcal{D}}

\paragraph{Problems}
\begin{enumerate}
\litem{(Uniform Noise Won't Help You.)}


One of the suggested ``solutions'' for a dataset with far more
positive labels for group $g_1$ than group $g_2$ suggested by ``Data
preprocessing techniques for classification without discrimination'',
\emph{massaging}, suggests flipping some fraction of the labels of the
datapoints in $g_1$ from negative to positive.

Suppose the current dataset is \emph{linearly separable} (namely,
there is some linear classifier with accuracy $1$ in training).

\begin{enumerate}
\item Describe a set of points in $\R^2$, broken into groups
  $g_1, g_2$:
  \begin{itemize}
    \item with equally sized groups $g_1, g_2$
    \item with 50\% positive and negative labels for group $g_1$
    \item  40\% positive labels for $g_2$
    \item which is linearly separable
    \end{itemize}

    Then, find a subset of $g_2$'s negative labels (large enough give
    the groups have equal number of positive labels if flipped) that,
    if swapped, would not change the set of accuracy-maximizing linear
    separators in training.
  \item Do the same, but rather than picking a carefully designed set
    of points in $g_2$, argue that a \emph{uniformly random} subset of
    negatively labeled points in $g_2$ (of the appropriate size) has
    low probability of changing the set of accuracy-maximizing linear
    classifiers. Give an estimate of this probability.

 
  \end{enumerate}

\litem{(Project Outline.)}

Please describe what you'd like to do for your course
project. Describe the motivation for this project, what concrete tasks
you plan to do, and what the final work product of these tasks will
be. Is there a particular context you're thinking about (e.g., hiring,
advertising) as you formulate this question?  If you're planning to do
an empirical project, list the dataset(s) you will use or gather.

It isn't necessary that you do precisely what you outline above for
the final project, but it should help you start thinking about your
work to answer these questions.
\end{enumerate}
\end{document}

%%% Local Variables:
%%% mode: latex
%%% TeX-master: t
%%% End:
